\section*{第 \ref{ch:g12:kinematics-2d} 章——二维运动学}

\begin{solution}{exo:g12:kinematics-2d:1}
$\vect v = (3.0, 4.0)$，恒定；速率 $\sqrt{3.0^2 + 4.0^2} =
\qty{5.0}{m/s}$。轨迹为直线 $y = \tfrac43 x$。
速度矢量恒定：匀速直线运动。
\end{solution}

\begin{solution}{exo:g12:kinematics-2d:2}
直线，间距相等：匀速直线运动；
$v = 0.35/0.10 = \qty{3.5}{m/s}$。若间距逐渐增大，
则说明速率在增大。
\end{solution}

\begin{solution}{exo:g12:kinematics-2d:3}
$v(t) = 4.0\,t$，$a = \qty{4.0}{m/s^2}$：加速度恒定，
匀加速直线运动。$v(3.0) = \qty{12}{m/s}$。
\end{solution}

\begin{solution}{exo:g12:kinematics-2d:4}
$f = 1/8.0 = \qty{0.125}{Hz}$;
$v = 2\pi \times 4.5/8.0 \approx \qty{3.5}{m/s}$;
$a = v^2/R = 3.53^2/4.5 \approx \qty{2.8}{m/s^2}$，指向
旋转木马的中心轴。
\end{solution}

\begin{solution}{exo:g12:kinematics-2d:5}
斜率：$a = (10.0 - 2.0)/4.0 = \qty{2.0}{m/s^2}$。面积（梯形）：
$d = \tfrac12 (2.0 + 10.0) \times 4.0 = \qty{24}{m}$.
\end{solution}

\begin{solution}{exo:g12:kinematics-2d:6}
$v = \qty{88.9}{m/s}$: $a_n = 88.9^2/6000 \approx \qty{1.3}{m/s^2}
\approx 0.13\,g$。由 $a_n = v^2/R$ 可知，$v$ 越大，所需 $R$ 越大，
才能使侧向加速度（以及轨道和乘客所受的力）保持在可接受范围内——
弯道倾斜改变的是由谁提供这一加速度，
而不改变其大小。
\end{solution}

\begin{solution}{exo:g12:kinematics-2d:7}
$v(t) = -2.4\,t + 8.0$；$a = \qty{-2.4}{m/s^2}$，加速度恒定。换向
在 $v = 0$ 时，$t = 8.0/2.4 \approx \qty{3.3}{s}$，
位于 $x \approx \qty{16.3}{m}$ 处。之前：沿正向，减速；之后：沿负向，
速率增大——全程匀加速。
\end{solution}

\begin{solution}{exo:g12:kinematics-2d:8}
$v_1 = 0.28/0.40 = \qty{0.70}{m/s}$;
$v_2 = (0.54 - 0.10)/0.40 = \qty{1.10}{m/s}$;
$v_3 = (0.88 - 0.28)/0.40 = \qty{1.50}{m/s}$。速度每
经过 \qty{0.20}{s} 即增加 \qty{0.40}{m/s}，$a = \qty{2.0}{m/s^2}$，加速度恒定
——匀加速直线运动。
\end{solution}

\begin{solution}{exo:g12:kinematics-2d:9}
$T = \qty{2.36e6}{s}$: $v = 2\pi \times \num{3.84e8}/\num{2.36e6}
\approx \qty{1.0e3}{m/s}$;
$a_n = v^2/R \approx \qty{2.7e-3}{m/s^2}$。比值：
$g/a_n \approx 3600 = 60^2$——万有引力随距离的
平方成反比减弱（见\cref{ch:g10:universal-gravitation}），正好
解释了月球的绕转。
\end{solution}

\begin{solution}{exo:g12:kinematics-2d:10}
$f = 1200/60 = \qty{20}{Hz}$, $T = \qty{0.050}{s}$;
$v = 2\pi R f = 2\pi \times 0.25 \times 20 \approx \qty{31}{m/s}$;
$a_n = v^2/R \approx \qty{3.9e3}{m/s^2} \approx 400\,g$。滚筒
壁提供衣物所需的 $a_n$；经小孔时，没有任何东西
把水向内推，因此水沿切线方向飞出。
\end{solution}

\begin{solution}{exo:g12:kinematics-2d:11}
$v_0 = \qty{25}{m/s}$: $t_s = 25/7.5 \approx \qty{3.3}{s}$;
$d = v_0^2/(2a) = 625/15 \approx \qty{42}{m}$。图像为
一条直线段，自 \qty{25}{m/s} 降至 $0$，时刻为 $t_s$。
\end{solution}

\begin{solution}{exo:g12:kinematics-2d:12}
$\vect v = (-R\omega\sin\omega t,\, R\omega\cos\omega t)$，其
模恒为 $R\omega$；它与 $\vect{OM}$ 的点积为零，
故与半径垂直，即沿切线方向。再次求导，
得 $\vect a = -\omega^2\vect{OM}$：指向 $O$，模为
$\omega^2 R = (v/R)^2 R = v^2/R$.
\end{solution}

\begin{solution}{exo:g12:kinematics-2d:13}
$a_t = \qty{2.5}{m/s^2}$（沿反方向），$a_n = 20^2/80 =
\qty{5.0}{m/s^2}$;
$a = \sqrt{2.5^2 + 5.0^2} \approx \qty{5.6}{m/s^2}$；与运动方向的
夹角：$\theta = 180^\circ - \arctan(5.0/2.5) \approx 117^\circ$
（即与正后方的夹角为 $63^\circ$）。$5.6 < 7.0$：在抓地力
范围内，余量不大。
\end{solution}

\begin{solution}{exo:g12:kinematics-2d:14}
起动阶段：$t_1 = 20/1.2 \approx \qty{17}{s}$，
$d_1 = 20^2/(2 \times 1.2) \approx \qty{167}{m}$。匀速行驶：
\qty{40}{s}，$\qty{800}{m}$。制动阶段：$t_3 = 20/1.5 \approx
\qty{13}{s}$，$d_3 = 400/3.0 \approx \qty{133}{m}$。总计：
$\qty{1.1e3}{m}$，用时 \qty{70}{s}；平均速率 $1100/70 \approx
\qty{16}{m/s}$（\qty{57}{km/h}）。$v(t)$ 图像为梯形，其面积
$\tfrac12(70 + 40) \times 20 = \qty{1100}{m}$，与总路程相符。
\end{solution}

\begin{solution}{exo:g12:kinematics-2d:15}
$v = 2\pi \times \num{4.22e7}/\num{86164} \approx \qty{3.1e3}{m/s}$;
$a_n = v^2/R \approx \qty{0.22}{m/s^2}$。地球在该距离处的
引力提供了它；周期与地球的自转周期相等，因此
卫星与地面同步转动，始终悬停在赤道上空的
同一点（见\cref{ch:g12:satellites-kepler}）。
\end{solution}

\begin{solution}{pb:g12:kinematics-2d:1}
\textbf{1.} 运动具有相对性（\cref{ch:g10:relative-motion}）：每个
位置和速度都必须指明其所属的参考系。此处取地面（道路）
参考系，$x$ 沿道路方向，从第一帧的位置量起。

\textbf{2.} 各时间间隔内：$9.0/0.50 = \qty{18.0}{m/s}$——四个
数值相等。

\textbf{3.} 直线道路上，相等时间内位移相等，
为匀速直线运动。

\textbf{4.} $v_0 = \qty{18.0}{m/s} = \qty{64.8}{km/h}$.

\textbf{5.} $\vect a = \vect 0$.

\textbf{6.} $v(t) = -a\,t + v_B$; $x(t) = -\tfrac12 a t^2 + v_B\,t$.

\textbf{7.} $v(t_s) = 0$ 得 $t_s = v_B/a$；则
$d_b = -\tfrac12 a (v_B/a)^2 + v_B^2/a = v_B^2/(2a)$.

\textbf{8.} $v_B = \sqrt{2 a d_b} = \sqrt{2 \times 6.0 \times 27} =
\sqrt{324} = \qty{18.0}{m/s}$——恰与拍摄到的车速相同：司机
刹车前从未减速。

\textbf{9.} 反应距离：$18.0 \times 1.0 = \qty{18}{m}$；总距离
$18 + 27 = \qty{45}{m}$.

\textbf{10.} $d_b = v_B^2/(2a) \propto v_B^2$：$v_B$ 加倍时
$d_b$ 变为原来的 $4$ 倍。当车速为 \qty{50}{km/h}（$\qty{13.9}{m/s}$）时：
$d_b = 13.9^2/12.0 \approx \qty{16}{m}$.

\textbf{11.} 速率不变，但速度的方向在改变：
$\vect v$ 不恒定，因此汽车做加速运动。

\textbf{12.} 转过四分之一圆周：$2\pi R/4 = \qty{19.6}{m}$；
$v = 19.6/2.5 \approx \qty{7.9}{m/s}$ (\qty{28}{km/h}).

\textbf{13.} $a_n = 7.85^2/12.5 \approx \qty{4.9}{m/s^2}$，方向指向
环岛中心。

\textbf{14.} $a_{\max} = \mu g = 0.70 \times 9.81 \approx
\qty{6.9}{m/s^2}$; $v_{\max} = \sqrt{\mu g R} =
\sqrt{6.87 \times 12.5} \approx \qty{9.3}{m/s}$ (\qty{33}{km/h}).
$7.9 < 9.3$：抓地力充足——与
干净、无打滑的车道相符。

\textbf{15.} $T = 2\pi R/v = 78.5/7.85 \approx \qty{10.0}{s};
f = \qty{0.10}{Hz}$.

\textbf{16.} $64.8 - 50 \approx \qty{15}{km/h}$ 超速——约
高出限速 $30\%$。

\textbf{17.} 制动作用距离为 $41 - 18 = \qty{23}{m}$：
$v^2 = v_B^2 - 2ad = 324 - 2 \times 6.0 \times 23 = 48$，因此
撞击时 $v \approx \qty{6.9}{m/s} \approx \qty{25}{km/h}$。

\textbf{18.} 当速度为 \qty{13.9}{m/s} 时：$13.9 \times 1.0 + 13.9^2/12.0 =
13.9 + 16.1 = \qty{30.0}{m} < \qty{41}{m}$——汽车在距货车
\qty{11}{m} 处停下。

\textbf{19.} $d(v) = v\,t_r + v^2/(2a)$：一项呈线性（反应
距离即 $v t_r$，在刹车生效前走过），一项呈平方（动能
部分 $v^2/(2a)$）——这就是为什么适度超速的代价
会不成比例地增大。

\textbf{20.} 行车记录仪与刹车痕迹相互吻合：车速 \qty{65}{km/h}，在
\qty{50}{km/h} 限速路段行驶；若按法定限速，汽车将在距货车
\qty{11}{m} 处停下，事故正是由多出的
\qty{15}{km/h} 引起的。
\end{solution}
